\section{教育背景}

\textbf{电子科技大学 (UESTC)} \hfill 2022年9月 --- 至今

\textbf{格拉斯哥大学 (UofG), 双学位项目} \hfill 2022年9月 --- 至今

\begin{itemize}
    \setlength\itemsep{-0.5em}
    \item \textbf{专业学位}: 电子与通信工程 (ECE) 学士; \MYhref{https://marcobisky.github.io/cv/score.pdf}{GPA: 3.87/4.0}, \MYhref{https://marcobisky.github.io/cv/rank.pdf}{排名: 2/164 (前1.2\%)}.
    \item \textbf{相关课程}: 信息论, 人工智能与机器学习, 随机过程, 数字电路设计等.
    \item \textbf{线上课程}: 抽象代数, 复分析, 微分几何, 凸优化等.
\end{itemize}


\section{研究与项目经历}

\textbf{边缘 TinyML 的 RISCV 加速器系统级软硬件协同设计} \hfill 2025年9月 --- 至今

\vspace{-0.5em} \begin{small} \textit{研究助理, \MYhref{https://scholar.google.com/citations?user=1NT1jFMAAAAJ&hl=en}{李耘 教授}, 电子科技大学} \end{small} \vspace{-0.5em}

\begin{itemize}
    \setlength\itemsep{-0.5em}
    \item 基于谷歌 2025 年最新 \MYhref{https://github.com/google-coral/coralnpu/tree/main}{Coral NPU} 端到端开源框架, 使用 RISCV Vector (RVV) 汇编和 C 内联函数移植并加速了 MobileViT 模型中的深度可分离卷积层 (DSC) 和注意力层, 实现了边缘设备上的图像识别推理.
    \item 使用 \texttt{Verilator}, \texttt{Cocotb}, \texttt{Bazel} 和 \texttt{CMakeLists} 进行仿真和构建, 在 Arty A100T FPGA 上测试实时摄像头输入、使用 UART 与 Linux 通信输出物体位置.
\end{itemize}

\MYhref{https://github.com/Marcobisky/YOPO}{\textbf{YOPO: You Pick Only Once --- 轻量级目标跟踪算法}} \hfill 2025年9月

\begin{itemize}
    \setlength\itemsep{-0.5em}
    \item 开发了一种只需一次初始选择的轻量级目标跟踪算法, 避免了每一帧上 DNN 前向传播的密集计算.
    \item 利用基于归一化互相关系数 (NCC) 的匹配和自适应核的更新, 实现了高精度跟踪颜色和大小渐变的目标.
\end{itemize}

\textbf{\MYhref{https://marcobisky.github.io/cv/drone.pdf}{自主四旋翼系统}的\MYhref{https://marcobisky.github.io/posts/quadrocopter-control/}{控制}与\MYhref{https://marcobisky.github.io/cv/drone-cv.pdf}{计算机视觉}} \hfill 2025年2月 --- 2025年6月

\begin{itemize}
    \setlength\itemsep{-0.5em}
    \item 开发了一款用于目标检测、路线规划和闭环飞行控制的自动四旋翼飞行器.
    \item 使用 \texttt{ROS2} 和 \texttt{OpenCV} 库实现了原创设计的计算机视觉算法, 用于实时着陆区域检测.
\end{itemize}

\MYhref{https://github.com/Marcobisky/my-riscv}{\textbf{完整单周期 RV32I CPU 内核的设计与可视化}} \hfill 2025年1月 --- 2025年3月

\begin{itemize}
    \setlength\itemsep{-0.5em}
    \item 使用 \texttt{Verilog} 从零手搓了一个单核、单周期 32 位 RISCV CPU, 使用 \texttt{Verilator} 仿真, 并在 \texttt{Digital} 软件中进行了工作原理可视化, 已在 \MYhref{https://github.com/Marcobisky/my-riscv}{Github} 上开源.
    \item 构建了完整的数据通路, 包括 PC、取指器、译码器、寄存器堆、ALU、基于 LRU 的 L1 缓存等, 支持 GPIO, IIC, UART 等基本外设.
    \item 使用 RISCV 汇编实现了一个引导程序, 以及使用 \texttt{C} 语言实现了基本的 delay 和 GPIO 库. 使用 RISCV GNU 工具链进行交叉编译和仿真.
\end{itemize}

\MYhref{https://github.com/Marcobisky/CNN-for-mbed}{\textbf{面向嵌入式系统的 CNN/LSTM}} \hfill 2024年2月 --- 2024年5月

\begin{itemize}
    \setlength\itemsep{-0.5em}
    \item 设计并基于 STM32 微控制器实现了 CNN 与 LSTM 模型的推理, 实现了高精度和低延迟的患者跌倒检测预警、体温监测及数据曲线的可视化.
    \item 本地化采集、标注并构建了三轴加速度时序数据集. 在 Linux 环境下完成模型训练, 使用 \texttt{C++} 和 MbedOS 库在 STM32 上进行了硬编码实现与算法加速.
\end{itemize}


\section{相关技能}

\begin{tabularx}{\linewidth}{@{}l X@{}}
\textbf{IT 技能} &  \normalsize{Latex, Quarto Markdown, Linux, \MYhref{https://www.bilibili.com/video/BV1AterevErt/?spm_id_from=333.1387.homepage.video_card.click&vd_source=42579e22289b6144ba0b2bdcf99834e3}{Manim}, \MYhref{https://github.com/Marcobisky}{Github}.}\\
\textbf{编程语言} &  \normalsize{\texttt{C/C++}, \texttt{Python}, \texttt{RISCV Assembly}, \texttt{Verilog}, \texttt{Makefile}, \texttt{Bazel}, \texttt{Chisel}, \texttt{Matlab}.}\\
\textbf{英语} &  \normalsize{英语四六级 (615, 573), \MYhref{[IELTS 7](https://marcobisky.github.io/cv/ielts-score.pdf}{雅思 7} (8/7/7/6.5)}\\
\end{tabularx}


\section{奖项}

\textbf{电子科技大学一等学业奖学金 (前 5\%)} \hfill 2023年12月, 2024年12月 \\
\textbf{国家奖学金 (前 0.2\%)} \hfill 2024年12月 \\
\textbf{\MYhref{http://www.moe.gov.cn/srcsite/A17/moe_794/moe_628/202408/t20240806_1144389.html}{全国第七届大学生艺术展演一等奖 (小提琴声部)}} \hfill 2024年9月

